\subsection*{3.5. Una resolución parcial plana canónica de un
$\mathbb Z$-módulo arbitrario}

Aplicaremos los resultados precedentes cuando $\underline C$ sea la
categoría de los grupos conmutativos del topos $\underline T$, es decir,
la categoría de los $\mathbb Z$-módulos de $\underline T$. Sea $P$ un
objeto de $\underline C$. Definiremos, funtorialmente en $P$, un complejo
de cadenas $L_\bullet(P)$ aumentado hacia $P$:

%%%%%%%%%%  scan idx 210  |  folio 199  %%%%%%%%%%

\begin{equation}\tag{3.5.1}
\begin{tikzcd}[column sep=small,row sep=large]
 0\arrow[r]&\overbrace{\mathbb Z[P\times P]\times
   \mathbb Z[P\times P\times P]}^{L_2}
   \arrow[r,"d_1"]
 &\overbrace{\mathbb Z[P\times P]}^{L_1}
   \arrow[r,"d_0"]
 &\overbrace{\mathbb Z[P]}^{L_0}
   \arrow[r]\arrow[d,"\varepsilon"']&0\\
 &&&P&{}
\end{tikzcd}\quad ,
\end{equation}
con $L_i(P)=0$ para $i\geq3$. Las componentes se muestran en (3.5.1),
con la notación de 1.4. Para describir los diferenciales $d_0,d_1$ y el
aumento $\varepsilon$, escribamos $[p]$ para la sección de
$\mathbb Z[P](S)$ definida por $p\in P(S)$, y análogamente $[p,q]$ y
$[p,q,r]$ para las secciones correspondientes de
$\mathbb Z[P\times P](S)$ y
$\mathbb Z[P\times P\times P](S)$. Pongamos entonces
\begin{equation}\tag{3.5.2}
\left\{
\begin{aligned}
 \varepsilon[p]&=p,\\
 d_0[p,q]&=[p+q]-[p]-[q],\\
 d_1[p,q]&=[p,q]-[q,p],\\
 d_1[p,q,r]&=[p+q,r]-[p,q+r]+[p,q]-[q,r].
\end{aligned}
\right.
\end{equation}
Estos homomorfismos de grupos conmutativos definen efectivamente un
complejo $L_\bullet(P)$ aumentado hacia $P$:
\[
 \varepsilon d_0=0,\qquad \sisi{d_1d_0}{d_0d_1}=0.
\]
La primera relación es inmediata. La segunda se descompone en dos
relaciones según los dos factores de $L_2(P)$; estas expresan,
respectivamente, la conmutatividad y la asociatividad de la ley de

%%%%%%%%%%  scan idx 211  |  folio 200  %%%%%%%%%%

composición sobre $P$.

El complejo aumentado depende funtorialmente de $P$. Un morfismo
$u:P\longrightarrow P'$ da un diagrama conmutativo
\[
\begin{tikzcd}[column sep=huge,row sep=large]
 L_\bullet(P)\arrow[r,"L_\bullet(u)"]\arrow[d]
 &L_\bullet(P')\arrow[d]\\
 P\arrow[r,"u"]&P'.
\end{tikzcd}
\]
Sus componentes son planas, como todo $\mathbb Z$-módulo libre.

\underline{Proposición} 3.5.3. \underline{Se tiene}
$H_1(L_\bullet(P))=0$, \underline{y el aumento} $\varepsilon$
\underline{induce un isomorfismo}
\[
 H_0(L_\bullet(P))\xrightarrow{\sim}P.
\]

Para la demostración puede hacerse la reducción usual al caso en que
$\underline T$ es la categoría de los conjuntos y $P$ es un
$\mathbb Z$-módulo ordinario, caso en el cual (3.5.1) será probablemente
conocido. Daremos, en cambio, una demostración directa que no usa esta
reducción, aplicando 3.4.7 al aumento
\[
 \varepsilon:L_\bullet(P)\longrightarrow P,
\]
donde $P$ se considera, como de costumbre, un complejo de cadenas
concentrado en grado cero. Queda por demostrar que, para todo objeto $G$
de $\underline C$, el funtor
\begin{equation}\tag{*}
 \varepsilon^*:\underline\Psi_P(G)\longrightarrow
 \underline\Psi_{L_\bullet(P)}(G)
\end{equation}
es una equivalencia de categorías.

%%%%%%%%%%  scan idx 212  |  folio 201  %%%%%%%%%%

La fuente de (*) es claramente isomorfa a la categoría de las extensiones
del $\mathbb Z$-módulo $P$ por $G$. Para describir su destino, se usa 1.4
para identificar las categorías de extensiones de $L_0,L_1,L_2$ por $G$,
respectivamente, con la categoría de los torsores bajo $G_P$, la categoría
de los torsores bajo $G_{P\times P}$ y la categoría de las parejas formadas
por un torsor bajo $G_{P\times P}$ y un torsor bajo
$G_{P\times P\times P}$.

Así, un objeto del destino de (*) es una pareja $(E,\varphi)$, donde $E$
es un torsor bajo $G_P$ y $\varphi$ es una trivialización del torsor $E'$
sobre $P\times P$ correspondiente a la imagen inversa por $d_0$ de la
extensión de $\mathbb Z[P]$ por $G$ definida por $E$. La trivialización
está sujeta a la compatibilidad con $\sisi{d_1d_0}{d_0d_1}=0$ descrita
más adelante. La fórmula (3.5.2) da inmediatamente
\[
 E'=\pi^*E\,\pr_1^*(E)^{-1}\,\pr_2^*(E)^{-1},
\]
y una trivialización de este torsor equivale a un isomorfismo, denotado
también por $\varphi$,
\[
 \varphi:\pr_1^*E\,\pr_2^*E\xrightarrow{\sim}\pi^*E.
\]
Al explicitar la compatibilidad de $\varphi$ con $d_0d_1=0$ se obtienen
la condición de conmutatividad (1.2.1) y la condición de asociatividad
(1.1.4.1), correspondientes a los dos factores de $L_2$. La afirmación
de que (*) es una equivalencia es, por tanto, precisamente el diccionario
de 1.1.6. Esto demuestra 3.5.3.

\underline{Observación} 3.5.4. a) Sería interesante reforzar 3.5.3
construyendo una resolución $L_\bullet(P)$ de $P$, funtorial en $P$, cuyas
componentes fueran funtorialmente isomorfas a sumas de
$\mathbb Z$-módulos libres engendrados por potencias cartesianas

%%%%%%%%%%  scan idx 213  |  folio 202  %%%%%%%%%%

de $P$. En vista de los isomorfismos de la forma (1.4.5), esto daría un
método para calcular parcialmente $\Ext^i(P,A)$ como término final de una
sucesión espectral cuyo término inicial se expresa mediante los grupos
\[
 H^j(P\times P\times\cdots\times P,A).
\]
Aquí solo obtenemos información muy parcial sobre $\Ext^1$, y ninguna
información sobre los $\Ext^i$ superiores. En todo caso, hay una manera
natural de continuar en un grado la resolución parcial
$L_\bullet(P)$.\footnote{L. Breen construyó una resolución de este tipo
(nota del redactor).}

b) Hay una variante de (3.5.1) y 3.5.3 en la categoría de los
$A$-módulos, donde $A$ es un anillo fijo de $\underline T$. La forma de
$L_\bullet(P)$ es más complicada porque debe ampliarse la descripción
1.1.6 (compárese con 2.10.3). Se obtiene
\begin{equation}\tag{3.5.5}
\left\{
\begin{aligned}
 L_0={}&A[P],\\
 L_1={}&A[P\times P]+A[A\times P],\\
 L_2={}&A[P\times P\times P]+A[P\times P]
       +A[A\times A\times P]+A[A\times P\times P].
\end{aligned}
\right.
\end{equation}
La estructura de complejo aumentado viene dada por (3.5.2), completada
con
\begin{equation}\tag{3.5.6}
\left\{
\begin{aligned}
 d_0[\lambda,p]={}&[\lambda p]-\lambda[p],\\
 d_1[\lambda,\mu,p]={}&[\lambda,\mu p]-[\lambda\mu,p]
                         +\lambda[\mu,p],\\
 d_1[\lambda,p,q]={}&[\lambda,p+q]-[\lambda,p]-[\lambda,q]
                       +\lambda[p,q]-[\lambda p,\lambda q].
\end{aligned}
\right.
\end{equation}
La relación $d_0d_1=0$ se sigue de la relación análoga sobre las
componentes que aparecen en (3.5.2), la cual, como vimos, equivale a la
asociatividad y la conmutatividad de la adición sobre $P$, junto con las
relaciones
\[
 d_0d_1[\lambda,\mu,p]=0,
 \qquad d_0d_1[\lambda,p,q]=0.
\]

%%%%%%%%%%  scan idx 214  |  folio 203  %%%%%%%%%%

Estas últimas expresan, respectivamente, la asociatividad de la acción
dada de $A$ sobre $P$,
\[
 \lambda(\mu p)=(\lambda\mu)p,
\]
y su distributividad respecto de la adición,
\[
 \lambda(p+q)=\lambda p+\lambda q.
\]

Para un torsor $E$ bajo $G_P$, una trivialización de $d_0^*E$ equivale a
una ley de composición interna sobre $E$ (que satisface las condiciones
de compatibilidad ya descritas entre la ley interna de $E$ y su estructura
de torsor), junto con una ley de composición externa de $A$ sobre $E$,
compatible con la acción de $A$ sobre $P$ y sobre $G$. La compatibilidad
de esta trivialización con la trivialización evidente de $d_1^*d_0^*E$
equivale a las condiciones ya descritas, que dicen que la ley interna de
$E$ es una ley de grupo conmutativo que extiende $P$ por $G$, junto con la
asociatividad de la ley externa y la distributividad de esta última
respecto de la ley interna. Estas condiciones dicen exactamente que $E$,
con sus dos leyes de composición y su estructura de torsor, es una
extensión de $A$-módulos de $P$ por $G$.

Se sigue, por esencialmente el mismo método que en 3.5.3, que
$L_\bullet$ es una resolución parcial del módulo $P$, es decir,
$H_1(L_\bullet)=0$. Sería de nuevo interesante prolongar esta resolución
parcial a una resolución infinita del mismo tipo, con componentes de la
forma
\[
 A[A\times A\times\cdots\times P\times P\times\cdots].
\]
Compárese con la cuestión análoga de a) para grupos conmutativos.

%%%%%%%%%%  scan idx 215  |  folio 204  %%%%%%%%%%
